\documentclass[11pt,a4paper]{article}
\usepackage[T1]{fontenc}
\usepackage[utf8]{inputenc}
\usepackage{lmodern}
\usepackage[margin=25mm]{geometry}
\usepackage{amsmath,amssymb,amsthm,mathtools}
\usepackage{booktabs,longtable,array}
\usepackage{microtype,xcolor}
\usepackage{enumitem,listings,hyperref}
\hypersetup{colorlinks=true,linkcolor=blue!45!black,citecolor=blue!45!black,
 urlcolor=blue!45!black,pdftitle={Explicit Diophantine equations for Con(ZFC)},
 pdfauthor={Research contribution prepared for ProveIt}}
\setlength{\emergencystretch}{2em}
\setlength{\parskip}{3pt}
\setlist{itemsep=2pt,topsep=5pt}
\newtheorem{theorem}{Theorem}[section]
\newtheorem{lemma}[theorem]{Lemma}
\newtheorem{proposition}[theorem]{Proposition}
\newtheorem{corollary}[theorem]{Corollary}
\theoremstyle{definition}
\newtheorem{definition}[theorem]{Definition}
\theoremstyle{remark}
\newtheorem{remark}[theorem]{Remark}
\newcommand{\Z}{\mathbb Z}
\newcommand{\N}{\mathbb N}
\newcommand{\Pos}{\mathbb Z_{>0}}
\newcommand{\Con}{\operatorname{Con}}
\newcommand{\pc}{\operatorname{pc}}
\newcommand{\val}{\nu_2}
\newcommand{\band}{\mathbin{\&}}
\newcommand{\src}[1]{\texttt{\detokenize{#1}}}
\newcommand{\pathref}[1]{\nolinkurl{#1}}
\newcommand{\Pbase}{P_{67}}
\newcommand{\Plow}{P_{68}}
\DeclareMathOperator{\Halt}{Halt}
\lstset{basicstyle=\ttfamily\footnotesize,breaklines=true,
 columns=fullflexible,keepspaces=true,frame=single}
\title{\textbf{Explicit Diophantine equations for
 \(\Con(\mathrm{ZFC})\)}\\[5pt]
 \large A 61-operation system and a 67-operation polynomial\\
 for the fixed blank computation K215}
\author{Research contribution prepared for the ProveIt repository}
\date{22 October 2026\\[3pt]
 \small Source snapshot \texttt{cdd582bec13f908df3711fa9e05d1ddac368a5db}}
\begin{document}
\maketitle
\begin{abstract}
We give a completely specified polynomial in fifteen strictly positive
integer unknowns whose solvability is equivalent, provably in ZFC, to the
existence of a finite ZFC refutation. Its negated solvability is therefore
equivalent to \(\Con(\mathrm{ZFC})\). The polynomial has exact degree \(155\)
and an explicit straight-line evaluation with \(67=39M+28A\) binary
arithmetic operations. An equivalent six-equation system uses
\(61=34M+27A\) operations. A one-operation tradeoff lowers the polynomial
degree to \(123\). A division-free substitution gives corresponding
equations in \(41\) unrestricted integer unknowns.

The construction specializes the repository's native-window/Pell compiler
to the fixed blank-input K215 machine, deleting the numerical input
interface. All four coefficient ports are bound to specific integers by
closed finite formulas and the pinned 430-row transition table. There is
no unspecified universal-machine index. The principal cost of the short
arithmetic description is coefficient height: the cell base is
\(B=2^{5^{67}}\). The coefficients and the approximately
\(1.42\cdot10^{21}\) local windows are defined exactly, but are not expanded.

This contribution consists of a mathematical proof, explicit arithmetic
circuits, source pins, and reproducible finite audits. The new Diophantine
bridge has not been checked by Lean or Rocq. The existing K215 endpoint
formalizations are identified separately. No global minimality claim and
no proof of \(\Con(\mathrm{ZFC})\) is made.
\end{abstract}
\clearpage
\tableofcontents
\newpage

\section{Result, domain, and the simplicity criterion}
The convention throughout the main construction is that all supplied
unknowns range over \(\Pos\). Every intermediate expression is an ordinary
integer polynomial; intermediate values need not be positive off the
zero set. Section~\ref{sec:integers} removes the positivity convention.

\begin{theorem}[Fixed consistency equation]\label{thm:main}
Let the four integers \(m,K,M_C,M_F^{\rm src}\) be those defined in
Sections~\ref{sec:machine}--\ref{sec:constants}, with the K215 transition
table at the source snapshot stated on the title page. Let \(\Pbase\) be
the polynomial in Section~\ref{sec:polynomial}. Then
\[
 \mathrm{ZFC}\ \vdash\
 \left[
  \exists \mathbf x\in\Pos^{15}\;\Pbase(\mathbf x)=0
  \ \longleftrightarrow\
  \exists p\;\operatorname{Proof}_{\mathrm{ZFC}}(p,\bot)
 \right].
 \tag{1}\label{eq:mainexist}
\]
In particular,
\[
 \mathrm{ZFC}\ \vdash\
 \left[
  \forall\mathbf x\in\Pos^{15}\;\Pbase(\mathbf x)\ne0
  \ \longleftrightarrow\ \Con(\mathrm{ZFC})
 \right].
 \tag{2}\label{eq:mainno}
\]
The same assertions hold for the six-equation system in
Section~\ref{sec:polynomial}, for \(\Plow\), and, after the explicit
substitution in Section~\ref{sec:integers}, over \(\Z^{41}\).
\end{theorem}
Here \(\operatorname{Proof}_{\mathrm{ZFC}}\) denotes an ordinary effective
finite proof predicate for ZFC. Changing among the standard effective
presentations of ZFC changes neither consistency nor the result: the
finite proof translations are formalizable in ZFC. The calibration to
the exact presentation used by the machine is explained in
Section~\ref{sec:formal}.

We minimize the size of a \emph{complete explicit arithmetic
description} within the repository's existing compiler framework.
Each binary addition, subtraction, or multiplication costs one operation,
including multiplication by a fixed numeral. Fixed numerals, reading
unknowns, reusing a register, and comparing two computed registers cost
zero. Thus the system and single-polynomial figures measure different
interfaces. A sum-of-squares row includes all residual subtractions,
squares, and additions.
\begin{center}\small
\begin{tabular}{lrrrr}
\toprule
Presentation & Unknowns & Equations & Operations & Degree\\
\midrule
Basic system & \(15\) positive & \(6\) & \(61=34M+27A\) & \(60\) maximum\\
Reduced-coefficient system & \(15\) positive & \(6\) & \(62=34M+28A\) & \(46\) maximum\\
\(\Pbase=0\) & \(15\) positive & \(1\) & \(67=39M+28A\) & \(155\)\\
\(\Plow=0\) & \(15\) positive & \(1\) & \(68=39M+29A\) & \(123\)\\
Basic sum of squares & \(15\) positive & \(1\) & \(78\) & \(120\)\\
Reduced sum of squares & \(15\) positive & \(1\) & \(79\) & \(92\)\\
\(\Pbase\), integer substitution & \(41\) integer & \(1\) & \(158\) & \(310\)\\
\(\Plow\), integer substitution & \(41\) integer & \(1\) & \(159\) & \(246\)\\
\bottomrule
\end{tabular}
\end{center}
The smallest complete arithmetic construction obtained here is the
61-operation system. These are upper bounds for displayed circuits,
not arithmetic-circuit lower bounds. Fifteen is not a minimum variable
count: classical universal Diophantine constructions can use nine
positive unknowns, with a very different degree and description-size
tradeoff \cite{jones}. No effective nine-variable specialization is
claimed to have been emitted in this archive.

\section{The polynomial and the six equations}\label{sec:polynomial}
Supply exactly these fifteen positive integers:
\[
 \mathbf x=(J,F,\alpha,z,f,h,i,T,s,w,\tau,\eta,\zeta,y,\gamma).
 \tag{3}
\]
The symbol \(K\) below is a fixed coefficient. The lower-case \(k\) is
a polynomial abbreviation, not a sixteenth unknown. Define
\begin{align}
 q&=mJ+1,& C&=q-F-\alpha,& X&=wq,&Y&=sq^3,& E&=XY,\notag\\
 k&=\eta+\zeta,& c&=kY+\eta,& a&=E+Y,\notag\\
 H&=4a+3,&\Delta&=a^2+H=(a+2)^2-1,&
 D&=X+ac+\gamma H. \tag{4}\label{eq:abbrev}
\end{align}
Next put
\begin{align}
 G&=(q-1)(q-F)+\alpha=q^2-C-qF,\notag\\
 R&=G(q^2-1)+(M_C+qM_F^{\rm src})J,\notag\\
 L_0&=EkY=XY^2k,\qquad
 V=c(Tf-1)-Rf^2,\qquad S=(i\Delta c^2)^2. \tag{5}\label{eq:pack}
\end{align}
Finally form the six factors
\begin{align}
 N_0&=\tau^2-L_0(L_0+k),\notag\\
 N_1&=D^2-\Delta c^2,\notag\\
 N_A&=S(V^2-y^2)+y^2,\notag\\
 N_K&=k-hE-R,\notag\\
 N_T&=(K+w)C+q-F-z(q-1),\notag\\
 N_S&=\Delta f^2-S. \tag{6}\label{eq:factors}
\end{align}
The requested single equation is
\[
 \boxed{\quad \Pbase=N_0N_1N_AN_KN_TN_S-\Delta=0.\quad}
 \tag{7}\label{eq:poly}
\]
Equivalently, one may use the system
\[
 \boxed{\quad N_0=N_1=N_A=N_K=N_T=1,\qquad N_S=\Delta.\quad}
 \tag{8}\label{eq:system}
\]
Abbreviations in (4)--(6) are to be substituted, not imposed as extra
equations with extra unknowns. All powers occurring in this polynomial
have fixed exponents. The very large exponents used later to define
the fixed coefficients are also fixed integers.

For the degree-\(123\) alternative, replace only \(N_A\) by
\[
 N_A'=\Delta(f^2-1)(V^2-y^2)+y^2,\qquad
 \Plow=N_0N_1N_A'N_KN_TN_S-\Delta.
 \tag{9}\label{eq:lowpoly}
\]
Section~\ref{sec:degrees} proves that this replacement preserves the
zero set; it is not an off-zero substitution of equal expressions.
The archive's JSON packets contain every gate, fixed-numeral port,
witness port, output register, and operation ledger.
Appendix~\ref{app:gates} prints the complete 67-gate source.

\section{The fixed machine and its local computation language}\label{sec:machine}
\subsection{Pinned K215 semantics}
The source table is
\pathref{Computability/BusyBeaver/Research/ZFC/candidate-k215/k215.tm}.
Its SHA-256 is
\begin{center}\small
\texttt{fdbcef0825886035c3091c8537a69535bc96f9e265f0ae45863543d40ae70090}.
\end{center}
There are 215 working states \(0,\ldots,214\), halt state \(215\),
two tape symbols \(0,1\), start state \(0\), and an initially blank
two-sided tape. The 430 rows specify a write, a left/right move, and
a next state. The first transition is
\[
 (0,0)\longmapsto(1,+1,93).
 \tag{10}
\]
The unique transition into the halt state is
\((68,1)\mapsto(1,+1,215)\).
For tableau padding only, complete the table by
\((215,b)\mapsto(b,0,215)\). This preserves the write and movement on
the transition entering the halt state. No prefix or auxiliary machine
is added.
The verifier independently parses the export, the literal Lean table
\src{KDense215.lean}, and the literal Rocq table \src{Table215.v}, and
compares all 430 rows. A finite bootstrap check is supplementary evidence;
it is not the unbounded halting theorem.

\subsection{A 439-symbol cell alphabet}
Let \(\mathcal A\) be the following alphabet, with the indicated integer
codes:
\begin{center}
\begin{tabular}{ll}
\toprule
Code & Meaning\\
\midrule
\(b\), \(b\in\{0,1\}\) & ordinary tape symbol \(b\), no head\\
\(2(q_0+1)+b\), \(0\le q_0\le215\) & symbol \(b\), head in state \(q_0\)\\
\(434,435,436\) & initialization symbols \(L,Q,R\)\\
\(437\) & vertical boundary \(V_{\rm bd}\)\\
\(438\) & horizontal boundary \(H_{\rm bd}\)\\
\bottomrule
\end{tabular}
\end{center}
The payloads of \(L,Q,R\) are respectively
\((0,\varnothing),(0,0),(0,\varnothing)\). Their phases are \(L,Q,R\).
Ordinary symbols have no phase. \(V_{\rm bd}\) has boundary bits
\((v,h_b)=(1,0)\), \(H_{\rm bd}\) has \((0,1)\), and all interior
symbols have \((0,0)\).
The exact local predicate is applied to a \(3\)-by-\(3\) window, in
row-major order, with time increasing from the top row to the bottom
row. It inspects its center, left, right, previous-time, and next-time
symbols; the four corners are unrestricted. Require:
\begin{enumerate}
\item The center's \(v\)-bit equals its next-time neighbor's \(v\)-bit.
The center and right neighbor are not both \(V_{\rm bd}\);
the center and next-time neighbor are not both \(H_{\rm bd}\).
No symmetric left/center or previous-time/center test is added
to this center predicate.
\item If the center and its right neighbor are both nonvertical-boundary
symbols, their \(h_b\)-bits agree.
\item At an interior center, a phase is present exactly when the
previous-time symbol is \(H_{\rm bd}\). At such a center, if the right
neighbor is nonvertical-boundary, the ordered pair of center/right
phases belongs to \(LL,LQ,QR,RR\). No left phase-pair test is added.
An initialization cell next to a left boundary is \(L\), and one
next to a right boundary is \(R\).
\item A head with \(H_{\rm bd}\) as its next-time neighbor must be in
state \(215\). A head must not move through an adjacent vertical boundary.
\item If the next-time center is not \(H_{\rm bd}\), its payload is
the deterministic three-neighbor update. A current center head writes
its tape symbol; otherwise that tape symbol stays. A new head can arrive
from the left moving right, the center staying, or the right moving
left. Multiple arrivals are forbidden. A vertical boundary is treated
as blank and headless. A horizontal-boundary neighbor cannot supply an
interior update.
\end{enumerate}
At a boundary center only the boundary conditions apply. The full,
executable, unambiguous predicate is \src{local_valid} in
\src{code/blank_helical.py}; the preceding clauses specify the same
finite relation.

\subsection{Start and an impossible distinguished window}
Two locally valid windows, in row-major code order, are
\[
 S_0=
 \begin{pmatrix}438&438&438\\434&435&436\\0&1&188\end{pmatrix},
 \qquad
 E_0=
 \begin{pmatrix}0&438&0\\0&437&438\\0&437&0\end{pmatrix}.
 \tag{11}\label{eq:SE}
\]
Every initialized \(Q\) in a globally valid tableau has window \(S_0\):
the actual first transition (10) forces the next row. Conversely,
\(S_0\)'s center is an initialized \(Q\).
The window \(E_0\) satisfies the center predicate, but cannot occur
globally. Its previous-time center is \(H_{\rm bd}\) with vertical bit
zero, while that symbol's next-time neighbor is \(V_{\rm bd}\) with
vertical bit one. The predicate at that \(H_{\rm bd}\) contradicts
vertical-bit propagation. We reserve the compiler's forbidden
selector \(1\) for this window. No genuine computation is removed
by forbidding it.

\begin{lemma}[Marked periodic tableau]\label{lem:tableau}
There exists a doubly periodic globally valid \(\mathcal A\)-tableau
containing \(S_0\) if and only if K215 halts on blank input.
For the forward direction, global uniqueness of \(S_0\) is unnecessary.
For the reverse direction, its width and height can be padded
independently to arbitrary sufficiently large integers, in particular
to sufficiently large powers of five.
\end{lemma}
\begin{proof}
Start with a periodic tableau containing \(S_0\).
On the initialization row containing its \(Q\), the phase graph
\(L\to L,Q;\ Q\to R;\ R\to R\) cannot form a cycle through \(Q\).
Horizontal periodicity therefore supplies vertical boundaries on both
sides. Choose the nearest ones. Vertical-bit propagation makes them
the sides of a finite strip. The horizontal boundary through the
previous-time neighbor extends across the strip, and temporal
periodicity supplies a next boundary. Choose the first one.
The first interior row is \(L^uQR^v\) with \(u,v\ge1\).
It has one head in state zero and all tape cells blank.
Induction gives the exact deterministic run in the strip; the head
cannot escape. Its last interior row has its unique head in state
215. Extending the strip by blank tape recovers the actual blank run.

Conversely take a finite halting run. Choose a width \(w_0=u+v+2\)
with positive margins containing every visited tape cell. Make
column zero vertical boundary, the first time row horizontal boundary
away from column zero, the next row \(L^uQR^v\), then the run followed
by enough absorbing halt rows. Call this rectangle \(G_0\), of height
\(H_0\). Width and height can be increased independently.
Extend it to the plane by
\[
 T_0(x,t)=G_0\bigl(x\bmod w_0,\,
                   (t+\lfloor x/w_0\rfloor)\bmod H_0\bigr).
 \tag{12}
\]
The vertical-boundary column permits the one-row time displacement
at each horizontal seam: its payload is absent, it is treated as
blank in updates, and horizontal boundary-bit propagation stops there.
All clauses remain valid. The periods are \((w_0,-1)\) and \((0,H_0)\).
The resulting canonical quotient has exactly one \(S_0\), and no
\(E_0\).
\end{proof}

\subsection{Window lift, ordering, and exact cardinality}
Let \(\mathcal W\subset\mathcal A^9\) be all center-valid windows.
Order them as
\[
 W_0=S_0,\quad W_1=E_0,\quad
 W_2,\ldots,W_{\kappa-1}
 \text{ in lexicographic row-major order after these exclusions}.
 \tag{13}
\]
This is a finite, unconditional definition. It performs no halting
test. The exact cardinality is
\[
 \kappa=38\,256\,607\,897\cdot439^4
       =1\,420\,903\,358\,357\,108\,792\,377.
 \tag{14}
\]
Appendix~\ref{app:counts} gives the cross count and its independent
finite check.
Use the two-column horizontal overlap and the two-row temporal overlap
as the ternary relation between a window and its two neighbors.
For any cyclic word with offsets \(0,-1,-h_0\), the overlaps identify
every entry of every window with the appropriate center symbol,
including when some offsets coincide. Thus they reconstruct an actual
periodic plane. Conversely actual windows satisfy the overlaps.
For the rectangle in the preceding proof, put \(N=w_0H_0\), let
\(x_S\) be the initial head column, and use
\[
 j=-(x-x_S)-w_0(t-1)\pmod N. \tag{15}
\]
The right and next-time offsets are \(-1,-w_0\), and \(S_0\) is at
index zero. This supplies precisely the cyclic interface of the
native compiler.

\section{Four completely specified coefficients}\label{sec:constants}
This section binds every fixed-numeral port in (4)--(7).
The different symbols separate compiler geometry from
the Pell variables: \(a_0=439\) and \(\kappa\) count symbols and
windows; neither is the polynomial's \(a\) or \(k\).

\subsection{Finite layout}
Let \(\ell\) be the least integer at least two such that
\(2^\ell>\kappa+1\), and set \(A_0=2^\ell\).
Define
\begin{align}
 p_0&=\max\left(0,
     \left\lceil\frac{\kappa-2\ell-21a_0-4}{2}\right\rceil\right),&
 j_{\max}&=9a_0+4+p_0,\notag\\
 \mu&=A_0-2+\sum_{j=1}^{j_{\max}}A_0^j,&
 m_0&=2\ell+30a_0+8+2p_0,\notag\\
 n_{\rm pos}&=m_0+1,&
 n_{\rm dum}&=n_{\rm pos}-\kappa-9a_0-4,\notag\\
 e_*&=\kappa+12a_0,&
 M_0&=\kappa+15a_0+n_{\rm dum}. \tag{16}
\end{align}
For \(r,c\in\{0,1,2\}\) and \(u\in\{0,\ldots,a_0-1\}\), put
\[
 e(r,c,u)=\kappa+3a_0+(2-r)3a_0+(2-c)a_0+u,\qquad
 j(r,c,u)=1+(3r+c)a_0+u. \tag{17}
\]
The finite set of allowed native positions is the disjoint union
\begin{align}
 \mathcal E={}&
 \{0,\ldots,\kappa-1\}\ \cup\ \{e(r,c,u)\}\notag\\
 &\cup\ \{e_*,\ldots,e_*+n_{\rm dum}-1\}
 \ \cup\ \{M_0,3M_0,9M_0,27M_0\}. \tag{18}
\end{align}
Its cardinality is \(n_{\rm pos}\). Define nonnegative integer clause
coefficients \(c_e\), for \(e\in\mathcal E\), as follows:
\begin{align}
 c_i&=1+\sum_{r,c=0}^2 A_0^{\,j(r,c,W_i[3r+c])}
          +\sum_{t=1}^4 A_0^{9a_0+t}
             &&(0\le i<\kappa),\notag\\
 c_{e(r,c,u)}&=A_0^{\,j(r,c,u)},\qquad
 c_{3^{t-1}M_0}=A_0^{9a_0+t}
             &&(1\le t\le4),\notag\\
 c_e&=0 &&\text{at a dummy position}. \tag{19}
\end{align}
Here \(W_i[3r+c]\) is an integer between zero and \(438\), determined
by (13).
Let
\begin{align}
 E_{\max}&=27M_0,\notag\\
 H_C&=E_{\max}+24M_0+3a_0+1,\notag\\
 T_1&=H_C+2E_{\max}+a_0+1,\notag\\
 T_2&=T_1+4E_{\max}+2. \tag{20}\label{eq:bands}
\end{align}
Let \(L\) be the least power of five at least \(T_2+2E_{\max}+2\).
For a positive integer \(v\), write
\(\operatorname{bits}(v)=\lfloor\log_2v\rfloor+1\). Put
\begin{align}
 U_0&=4(n_{\rm pos}+2)\bigl(28(\kappa+9a_0+4)+6\bigr)+8,\notag\\
 b_{\min}&=\max\bigl\{\ell(j_{\max}+1)+3,\,
                \ell(9a_0+4)+\operatorname{bits}(U_0)+1,\,
                4\bigr\}. \tag{21}
\end{align}
Let \(b\) be the least power of five at least \(b_{\min}\), and define
\[
 \rho=2^b,\qquad d=bL,\qquad B=\rho^L=2^d. \tag{22}
\]
These definitions contain only specified finite integer operations,
finite sums, and bounded finite enumeration.

\subsection{The masks and transport coefficients}
Define
\begin{align}
 D_C={}&\rho^{3a_0}+\rho^{H_C+a_0}
          +\rho^{8M_0}+\rho^{24M_0}
          +\sum_{e\in\mathcal E}c_e
              \bigl(\rho^{T_1-e}+\rho^{T_2-e}\bigr),\notag\\
 D_R={}&\rho^{H_C}, \tag{23}\label{eq:DC}
\end{align}
and
\begin{align}
 M_F^0={}&\mu(\rho^{T_1}+\rho^{T_2})
  +\sum_{\substack{0\le r\le1\\0\le c\le2\\0\le u<a_0}}
                    \rho^{e(r,c,u)}\notag\\
 &+\sum_{\substack{0\le r\le2\\0\le c\le1\\0\le u<a_0}}
                    \rho^{H_C+e(r,c,u)}
  +\rho^{9M_0}+\rho^{27M_0},\notag\\
 M_C={}&B-1-\sum_{\substack{e\in\mathcal E\\e\ne1}}\rho^e
                     -2\rho^{e_*},\qquad
 M_F=M_F^0+4. \tag{24}\label{eq:masks}
\end{align}
The four coefficients of the delivered polynomial are exactly
\[
 \boxed{\quad m=B-1,\qquad K=D_C+BD_R,\qquad
             M_C,\qquad M_F^{\rm src}=M_F+B-1.\quad} \tag{25}
\]
In particular \(M_F^{\rm src}\) is not the one-cell mask \(M_F\).
Applying \(M_F<B-1\) directly to \(M_F^{\rm src}\) would be incorrect.

\subsection{Concrete scalar values and coefficient size}
The layout formulas yield
\begin{align*}
 \ell&=71,&n_{\rm dum}&=1,\\
 p_0&=710451679178554391506,&
 j_{\max}&=710451679178554395461,\\
 m_0&=1420903358357108796332,&
 M_0&=1420903358357108798963,\\
 E_{\max}&=38364390675641937572001,&
 H_C&=72466071276212548748431,\\
 T_1&=149194852627496423892873,&
 T_2&=302652415330064174180879.
\end{align*}
Moreover
\[
 b=5^{33},\qquad L=5^{34},\qquad
 d=5^{67}
 =67762635780344027125465800054371356964111328125.
 \tag{26}
\]
Thus \(m=B-1\) alone has exactly \(5^{67}\) binary digits.
The archive specifies the exact integers by the finite formulas above;
it does not contain their infeasible binary or decimal expansions.
It also does not enumerate all \(\kappa\) windows. This is a short
coefficient \emph{description}, not a small coefficient-height result.

\subsection{Compiler invariants and the fixed modulo-five unit}
The clause padding gives
\(\pc(\mu)=\ell+j_{\max}-1\) and
\(m_0=2\pc(\mu)+12a_0+2\).
Consequently the masks obey
\begin{gather}
 0<M_C,M_F<B-1,\qquad
 M_C\equiv2\pmod4,\qquad M_F\equiv4\pmod8,\notag\\
 \pc(M_C)=d-m_0-1,\qquad
 \pc(M_F)=m_0+1,\qquad
 \pc(M_C)+\pc(M_F)=d. \tag{27}\label{eq:maskproperties}
\end{gather}
The unused high radix positions make the range inequalities strict.
The low selector zero is permitted; selector one is forbidden.
The first dummy position has digits \(0,1,2,3\) available, and the other
native positions have Boolean digits.
The conservative bounds
\[
 \sum_{e\in\mathcal E}c_e
    \le14(\kappa+9a_0+4)A_0^{9a_0+4},
 \qquad \mu<A_0^{j_{\max}+1}
 \tag{28}
\]
and (21) imply
\[
 \rho\ge
 \max\left\{
 4(n_{\rm pos}+2)\left(2\sum_e c_e+6\right)+8,\,
 2\mu+4,\,16\right\}. \tag{29}\label{eq:mass}
\]
This is the no-carry bound used in the native compiler.
The change in the second band in (20) removes a conditional coefficient
repair from the older construction. Since \(b\) is a power of five,
\(\rho\equiv2\pmod5\), and
\(T_2-T_1\equiv2\pmod4\). Each paired term in the sum in (23)
therefore cancels modulo five. Since \(a_0=439\equiv3\pmod4\),
\[
 2D_C-D_R
 \equiv2\bigl(2+3\rho^{H_C}+1+1\bigr)-\rho^{H_C}
 \equiv3\pmod5. \tag{30}\label{eq:unit5}
\]
No uncomputed test or optional high monomial remains in the recipe.
The larger band spacing preserves
\(T_2-T_1>2E_{\max}\), \(L-(T_2-T_1)>E_{\max}\), and
\(T_2+E_{\max}<L\), the support inequalities used by the compiler.

\section{Arithmetic soundness: every zero recovers a computation}
\subsection{Unit factors, including the exceptional boundary}
Write \(A=a+2\). Positive inputs give
\(\Delta=A^2-1>0\) before any equation is imposed.
Since \(S=\Delta^2i^2c^4\),
\[
 N_S=\Delta N_s,\qquad N_s=f^2-\Delta i^2c^4. \tag{31}
\]
Equation (7) is consequently equivalent to
\[
 N_0N_1N_AN_KN_TN_s=1. \tag{32}\label{eq:unitprod}
\]
Every factor is an integer unit.
The norm \(N_0=-1\) is impossible. Indeed, for any integer \(u\ge2\),
\[
 x^2-u(u+1)v^2=-1,\quad v>0 \tag{33}
\]
has no integer solution. Choose a minimal \(v\); then
\(uv<|x|<\sqrt{u(u+1)}v\), and multiplication by the norm-one
unit \(2u+1-2\sqrt{u(u+1)}\) produces a solution with
\[
 0<v'=(2u+1)v-2|x|<v. \tag{34}
\]
This is a contradiction. Apply this with \(u=XY^2\).
Also \(\Delta\equiv0\) or \(3\pmod4\), so neither \(N_1\) nor \(N_s\)
can be \(-1\). Since \(S\) is a square, \(N_A\) is congruent modulo
four to \(y^2\) or \(V^2\), and cannot be \(-1\). Thus
\[
 N_0=N_1=N_A=N_s=1,\qquad
 N_K=\varepsilon,\quad N_T=\nu,\quad
 \varepsilon,\nu\in\{-1,1\},\quad\varepsilon\nu=1. \tag{35}
\]
It is essential not to assume \(C>0\) yet. Since
\[
 q-F-z(q-1)=1-F-(z-1)(q-1)\le0 \tag{36}
\]
and \(K+w\ge2\), either transport unit forces \(C\ge0\).
The boundary \(C=0,F=2,z=1,\alpha=q-2\) does give \(N_T=-1\);
it cannot be discarded by the transport factor alone.
Nevertheless \(q-F=C+\alpha>0\), and
\begin{gather}
 q\le G=q\alpha+(q-1)C\le q^2-q,\notag\\
 0<(M_C+qM_F^{\rm src})J<2q^2-q-1. \tag{37}
\end{gather}
Using \(q\ge B\ge16\), we obtain
\begin{gather}
 q(q^2-1)<R<q^4-q^3+q^2-1<q^4-2,\notag\\
 R-2>3q+1,\qquad E\ge q^4>R+2,\qquad a=E+Y>R+2. \tag{38}\label{eq:weakbounds}
\end{gather}
These bounds are weaker than a bound sometimes used in the older
variable-input proof. They suffice here because the recovered auxiliary
target will be exactly \(R\).

\subsection{The deleted auxiliary coordinates are positive}
The strong norm gives
\[
 f^2=1+\Delta i^2c^4,\qquad
 S=\Delta(f^2-1),\qquad f>c^2>2c,\qquad R<a<\Delta. \tag{39}
\]
Therefore
\[
 S-(Rf^2+c)
  =(\Delta-R)f^2-\Delta-c
  \ge1+\Delta(i^2c^4-1)-c>0. \tag{40}
\]
The auxiliary equation is
\[
 SV^2-(S-1)y^2=1. \tag{41}
\]
It excludes \(V=0\). For \(v=|V|>1\), it implies
\(y\ge v+1\) and
\(v^2-1=(S-1)(y^2-v^2)\ge(S-1)(2v+1)\), hence
\(v\ge2S-1\). But \(T>0\) gives
\(V>-Rf^2-c>-S\), excluding negative \(V\) of that size.
The cases \(V=\pm1\) are excluded by \(V\equiv-c\pmod f\) and
\(f>2c,\ c>2\). Hence \(V>0\).
Restore, as mathematical expressions only,
\[
 o=cT-Rf,\qquad
 j=Tf-R\Delta i^2c^3-1. \tag{42}\label{eq:restore}
\]
They are integers and obey
\[
 of=V+c>0,\qquad jc=V+R>0,\qquad V=of-c=jc-R. \tag{43}
\]
Thus \(o,j>0\), before any rank theorem or computation decoding.
No division or extra witness has been inserted in the polynomial.

\subsection{Pell rank and the remaining signs}\label{sec:rank}
For an integer \(U>1\), define the standard Pell sequences by
\[
 (U+\sqrt{U^2-1})^n
   =\chi_U(n)+\psi_U(n)\sqrt{U^2-1},
 \qquad n\ge0. \tag{44}
\]
We use the established Pell facts from the repository's
generic strong-rank, signed step-down, and half-parameter proofs
\cite{pellrank,pellsigned,halfparameter}. These are mathematical
dependencies, supplied in the archive; they are not consequences
of the finite program checks.
The first norm, with \(P_0=2XY^2+1\), has fundamental unit
\((P_0,2)\), and consequently
\[
 \tau=\chi_{P_0}(n),\quad k=2\psi_{P_0}(n),\quad n\ge1. \tag{45}
\]
The positive main root similarly gives
\[
 D=\chi_A(p),\qquad c=\psi_A(p),\qquad p\ge1. \tag{46}
\]
The strong-rank lemma says that, if
\(c>A\Delta^2\) and a positive ordinary coefficient \(i_0\)
satisfies \((i_0c^2)^2=\Delta(f^2-1)\), then
\[
 f=\chi_A(m_1),\qquad c\mid m_1,\qquad m_1\ge c. \tag{47}
\]
The signed auxiliary step-down lemma, applied with \(m_1>2p\),
the auxiliary norm (41), \(V=of-c=jc-R>0\), and
\(0<p,R<c/2\), gives \(R\equiv\pm p\pmod c\), hence \(p=R\).
In this application its auxiliary coefficient is
\(S=(i_0c^2)^2\). The full congruence argument and its hypotheses
are recorded in \src{proofs/pell_lemma_audit.md}.

Here are the checks of all the size hypotheses.
Since \(P_0\equiv1\pmod E\), the first recurrence gives
\[
 2n\equiv R+\varepsilon\pmod E,\qquad
 n\ge(R-1)/2\ge24. \tag{48}
\]
The inequalities \(P_0>A\) and \(c>k\) imply \(p>n\) by
Pell monotonicity. Standard growth then gives
\(c>A\Delta^2\) and \(c>2p\). Also
\(k=R+\varepsilon+hE>2R\), so \(c>k>2R\).
Use \(i_0=\Delta i\) in (47); its premise follows from (39).
Then \(m_1\ge c>2p\). Equations (41)--(43) verify the auxiliary
hypotheses. We conclude
\[
 p=R. \tag{49}
\]
Write \(2n=R+\varepsilon+vE\). The least positive residue in (48)
gives \(v\ge0\). If \(v\ge1\), (38) gives \(2n>2R\),
contradicting \(n<p=R\). Hence \(2n=R+\varepsilon\).
If \(\varepsilon=-1\), then \(p=2n+1\).
Let \(Q_0=\chi_A(2)=2A^2-1>P_0\). Pell duplication yields
\[
 \psi_A(2n)=2A\psi_{Q_0}(n)\ge2A\psi_{P_0}(n)=Ak. \tag{50}
\]
Since \(A>Y+1\), this would give
\(c=\psi_A(2n+1)>k(Y+1)\). It contradicts
\[
 kY<c<k(Y+1), \tag{51}
\]
which follows directly from the two positive slacks \(\eta,\zeta\).
Therefore \(\varepsilon=\nu=1\). The transport factor now excludes
\(C=0\), so \(C>0\).
This proves the equivalence of (7) and the individual equations (8).

\subsection{Power extraction before mask decoding}
We explain why the asymmetric scale \(X=wq\) suffices.
Let \(r_0=(R-1)/2\), now an integer, and put
\(\xi=(X+1)^{2r_0}/X^{r_0}\).
The standard ratio estimates for (45)--(46), using
\(P_0>A,\ 6XY^2>a,\ X\ge q\ge16,\ Y\ge q^3\), give
\[
 \xi<\frac{2c}{k}<2Y+2,\qquad
 Y>X^{r_0}/3,\qquad a>X^{r_0+1}/3. \tag{52}
\]
These lower estimates do not require \(q^3\mid X\) or the older
stronger upper bound on \(R\); see \cite{asymmetric}, Section 3.
In particular \(X<a\) and
\[
 2^R=2\cdot4^{r_0}<X^{r_0+1}/3<a. \tag{53}
\]
The Pell recurrence for \(\chi_A(t)-a\psi_A(t)\) is congruent
modulo \(H=4a+3\) to \(2^t\). Thus \(D=X+ac+\gamma H\)
implies \(X\equiv2^R\pmod H\). Both positive representatives
are below \(a<H\), proving
\[
 X=2^R. \tag{54}
\]
Because \(q\mid X\), \(q=2^t\) for an integer \(t>0\).
The bound \(R>3q\) now implies \(q^3\mid X\).
Thus one may restore the symmetric scale \(\widetilde w=X/q^3\)
in the proof and apply the full half-binomial kernel
\cite{halfbinomial}. Its hypotheses
\(q\ge16,\ 3q+1\le R<q^4\) hold by (38).
It gives
\begin{gather}
 Y=\frac12\left\lfloor
             \frac{(X+1)^{R-1}}{X^{(R-1)/2}}\right\rfloor,\notag\\
 \val(Y)=\pc(R)-2. \tag{55}
\end{gather}
This uses the central-binomial valuation
\(\nu_2\binom{2r_0}{r_0}=\pc(r_0)\);
all the remaining terms in the binomial floor are divisible by \(X\).
As \(q^3\mid Y\), (55) gives \(\pc(R)\ge3t+2\).
Finally \(B-1=2^d-1\) divides \(q-1=2^t-1\).
The order of two modulo \(2^d-1\) is \(d\), so \(t=dN\)
for some \(N\ge1\). Thus
\[
 q=B^N,\qquad J=1+B+\cdots+B^{N-1}. \tag{56}
\]

\subsection{Inverse packing and synchronization}
Let
\[
 S'=(C-1)+qF,\qquad
 T'=(M_CJ+1)+q(M_FJ-1). \tag{57}
\]
The source's shifted mask coefficient gives the exact identity
\[
 R=(q^2-S')(q^2-1)+T'. \tag{58}
\]
The now established \(C,F>0\) and \(C+F<q\) imply \(0<S'<q^2\).
Equation (27) and the repunit (56) give
\[
 0<T'<q^2-1,\qquad\pc(T')=dN+2. \tag{59}
\]
Adding one to the low block adds one population bit; subtracting
one from the high block with exactly two trailing zeros also adds
one bit.
The elementary inverse-population lemma states that for
\(\Lambda=2^n,\ 1\le S_1\le\Lambda,\ 0<T_1'<\Lambda-1\),
\[
 \pc\bigl((\Lambda-S_1)(\Lambda-1)+T_1'\bigr)\le n+\pc(T_1'),
 \tag{60}
\]
with equality exactly when \(S_1<\Lambda\) and
\(S_1\band T_1'=0\).
Splitting into its two \(n\)-bit blocks proves the assertion:
when \(S_1+T_1'<\Lambda\), equality means no addition carry; when
\(S_1+T_1'\ge\Lambda\), the outgoing carry and discarded leading bit
make the inequality strict. The boundary \(S_1=\Lambda\) is also
strict. A detailed proof is included in \cite{compiler75}.
Apply (60) to (58). The lower bound \(\pc(R)\ge3dN+2\)
forces equality, and therefore
\[
 (C-1)\band(M_CJ+1)=0,\qquad
 F\band(M_FJ-1)=0. \tag{61}\label{eq:typed}
\]
The first condition restores Start at the origin with no carry,
forbids selector one in every cell, and restricts native digits to
the positions in (18). It permits additional Start occurrences;
their exclusion is unnecessary.

The remaining compiler argument is the native coefficient proof
\cite{compiler76,compiler75}, with the layout and masks specified
above. We record its relevant content to make the change of local
language explicit. The no-carry bound (29) keeps every unshifted
coefficient at most \(\rho/4-2\).
In a binary rotation by an arbitrary exponent \(r\), write
\(r=bt_0+\epsilon_0\), \(0\le\epsilon_0<b\).
The extra upper bit at the dummy position means that the possible
rotated support is
\((\mathcal E\cup\{e_*+1\})+t_0\) modulo \(L\).
For \(\epsilon_0\le b-2\), a rotated digit is at most \(3\rho/4\);
for \(\epsilon_0=b-1\), it is at most \(\rho/2+1\).
The sum with the unshifted field still has no carry.
Let \(\widehat F\) be the sum of the two native convolutions
and that cyclic temporal rotation, represented in the \(N\)-cell
word. These bounds give \(0<\widehat F<q-1\).
The already proved equation \(N_T=1\), together with
\(X=wq\equiv w\pmod{q-1}\), gives
\(F\equiv\widehat F\pmod{q-1}\).
Since \(0<F<q\), the two representatives are equal:
\(F=\widehat F\). Thus the field-mask tests in (61) apply to the
actual carry-free field before any clause decoding is used.

The two distant clause bands cannot both meet that rotated support.
An uncontaminated band enforces the selector/copy/anchor clauses:
at occupied cells it specifies one window, its matching
payload-copy bits, and its four anchors. At the marked origin the
anchor tests require two distinct rotated digits to be odd.
For \(1\le\epsilon_0\le b-2\) all rotated digits are even.
For \(\epsilon_0=b-1\) only the single upper-dummy spill position
can be odd. Both cases are impossible.
Thus \(\epsilon_0=0\). The anchor pair differences then force
\(t_0=0\pmod L\). The rotation is a whole-cell shift \(B^{h_0}\).
The native occupancy propagation from Start and the copy tests
enforce the horizontal and temporal window overlaps.
The temporal overlap of Start's horizontal-boundary row and its
initialization row excludes \(h_0=0\pmod N\).
No unique-End or unique-Start theorem is used here.

The transport factor used here really gives the required congruence:
\[
 (K+w)C-F=(z-1)(q-1),\qquad X=wq\equiv w\pmod{q-1}. \tag{62}
\]
If a positive quotient for the unsheared form is wanted,
\[
 (K+X)C=F+(z-1+wC)(q-1),\qquad z-1+wC>0. \tag{63}
\]
Therefore the supplied \(F\) is the actual carry-free field,
and (61) decodes a marked periodic tableau for the local language
of Section~\ref{sec:machine}. The cyclic word pulls back along
\((x,t)\mapsto-x-h_0t\pmod N\), with periods \((N,0),(0,N)\);
neither \(h_0\mid N\) nor coprimality is needed.
Lemma~\ref{lem:tableau} proves that K215 halts. This completes
arithmetic soundness for every positive zero, including the
initially possible \(C=0\) branch.

\section{Completeness at the actual packed index}\label{sec:complete}
Suppose K215 has a finite blank halting run. We construct positive
unknowns without assuming any bound on the length of that run.
By Lemma~\ref{lem:tableau}, choose independent sufficiently large
powers of five \(w_0,H_0\), with \(H_0\ge25\),
\[
 N=w_0H_0>25d,\qquad q=B^N. \tag{64}
\]
Use the canonical cyclic word with one Start, no forbidden window,
and all dummy bits initially zero. Let \(C_0\) be its native word
and \(F_0\) its actual field for temporal shift \(B^{w_0}\).
Let \(R_0\) be the packed rank (5) evaluated at these \(C_0,F_0\).
The source and mask proofs in \cite{compiler75} show that (61)
holds. In particular the origin's extra field tests hold: the
temporal predecessor of the unique Start is not Start.

We must tune the \emph{actual value of \(R\)}, rather than choose a
formal exponent independently. Turn on ordinary low dummy bits in
cells from
\[
 \{4j:0\le j<N/5\}\cup\{1\}. \tag{65}
\]
These cells do not wrap under either shift \(1\) or \(w_0\).
For a selected cell \(v\), the exact changes are
\begin{align}
 \delta C&=\rho^{e_*}B^v,\notag\\
 \delta F&=(D_C+BD_R+B^{w_0})\rho^{e_*}B^v,\qquad
 \delta R=-\Gamma B^v, \tag{66}
\end{align}
where
\[
 \Gamma=\rho^{e_*}(q^2-1)
         \bigl[1+q(D_C+BD_R+B^{w_0})\bigr]. \tag{67}
\]
Because \(b,L,N,w_0\) are powers of five,
\(\rho,B,q,B^{w_0}\equiv2\pmod5\).
The bracket in (67) is \(2D_C-D_R\equiv3\pmod5\);
hence \(\Gamma\) is a unit modulo \(dN\).

For completeness, the Boolean subset lemma \cite{fiveadic} says
that every residue modulo \(dN\) is a sum of a subset of the
distinct weights \(B^v\) with \(v\) from (65).
Indeed \(B^4\) has order \(N/5\)
modulo \(dN\) and its powers enumerate exactly the residues
\(1+5dt\). For a desired residue, optionally subtract the weight
\(B\) to make its residue modulo \(5\) nonzero. Its residue
\(k_0\) modulo \(5d\) then satisfies \(1\le k_0<5d<N/5\)
and is a unit modulo \(N/5\). A cyclic interval of \(k_0\)
distinct residues \(t\) can be chosen to have any desired sum,
since its sum is
\(k_0t_0+k_0(k_0-1)/2\). This gives a genuine Boolean subset.
The order assertion follows from
\(\nu_5(2^{4dj}-1)=\nu_5(d)+1+\nu_5(j)\).
Apply this lemma to the target
\((R_0-dw_0)\Gamma^{-1}\pmod{dN}\).
The modified actual word and actual field give
\[
 R\equiv dw_0\pmod{dN}. \tag{68}
\]
The dummy bits are ignored by all clauses. Their upper bits remain
zero, all masks remain valid, and the canonical pointwise digit
bound gives
\[
 2C+F\le(\rho/4+1)\frac{q-1}{\rho-1}<q/3. \tag{69}
\]
Thus \(J=(q-1)/(B-1)>0\) and
\(\alpha=q-F-C>0\). This verifies the stronger slack built into
our definition of \(C\), including after the dummy adjustment.
By the exact shifted mask identity,
\[
 3q+1\le R<q^4,\qquad R\equiv3\pmod4,\qquad
 \pc(R)=3dN+2. \tag{70}
\]
For the congruence, (61) gives \(C\equiv1\pmod4\);
the remainder \(T'\) in (58) is \(3\pmod4\), and
\(S'\equiv0\pmod4\), so (58) gives \(R\equiv3\pmod4\).

Now invoke the positive half-binomial converse at this \emph{fresh,
actual} \(R\). Set \(r_0=(R-1)/2\), \(X=2^R\), and
\[
 Y=\frac12\left\lfloor
          \frac{(X+1)^{2r_0}}{X^{r_0}}\right\rfloor. \tag{71}
\]
Equation (70) and the exact valuation imply
\(q^3\mid Y\); also \(q\mid X\). Let \(s=Y/q^3\) and
\(w=X/q\). Both are positive integers.
Set \(a=Y(X+1)\), \(A=a+2\), \(\Delta=A^2-1\),
\(E=XY\), \(P_0=2XY^2+1\), and
\begin{align}
 c&=\psi_A(R),& D&=\chi_A(R),&
 k&=2\psi_{P_0}(r_0+1),& \tau&=\chi_{P_0}(r_0+1),\notag\\
 \eta&=c-kY,&\zeta&=k-\eta,&
 h&=(k-R-1)/E,&
 \gamma&=(D-X-ac)/(4a+3). \tag{72}
\end{align}
The kernel's strict ratio bounds and congruences make all these
integers positive. The supplied quotient \(h\) here is different
from the external helical width \(w_0\).
For the normalized strong auxiliary coordinates, let
\[
 m_1=2cR,\qquad
 f=\chi_A(m_1),\qquad i=\psi_A(m_1)/c^2,\qquad
 \rho_1=\Delta\psi_A(m_1). \tag{73}
\]
The quotient \(i\) is integral: the square-root coefficient of
\((\chi_A(R)+c\sqrt\Delta)^{2c}\) is divisible by \(c^2\);
its linear square-root term contains \(2c^2\), and all higher
odd terms contain at least \(c^3\).
The canonical minus construction \cite{halfbinomial,normalized}
at \(R\equiv3\pmod4\) gives
\[
 y=\psi_{\rho_1}(R),\quad
 V=\chi_{\rho_1}(R)/\rho_1,\quad
 o=(V+c)/f,\quad j=(V+R)/c \tag{74}
\]
as positive integers satisfying the required norms.
Since \(of+R=c(j+1)\) and \(f^2\equiv1\pmod c\),
\(c\mid o+Rf\). Therefore
\[
 T=(o+Rf)/c>0 \tag{75}
\]
is integral and reconstructs exactly \(V=c(Tf-1)-Rf^2\).
Finally (68) gives \(X\equiv B^{w_0}\pmod{q-1}\), and the
same holds for \(w=X/q\). The actual field therefore makes
\((K+w)C-F\) divisible by \(q-1\).
In fact \(w=2^R/q>q>F\), by (70), so its quotient is
positive. Put
\[
 z=1+\frac{(K+w)C-F}{q-1}>0. \tag{76}
\]
Every equation in (8) now holds. Divisions and exponential formulas
in this section construct existential witnesses in the proof; they
are not operations in the delivered polynomial.

\section{Why the equivalence is provable in ZFC}\label{sec:formal}
The pinned Lean endpoint \cite{kendpoint} states, for the exact
machine \src{machine215},
\begin{lstlisting}
blank_halts_iff_zfc_refutation :
  BlankHalts machine215 <-> BProv ZFCax_s [] Form.fBot

nonhalting_iff_zfc_consistency :
  (not (BlankHalts machine215)) <->
  (not (BProv ZFCax_s [] Form.fBot))
\end{lstlisting}
The Rocq endpoint \src{CorrectDense215.v} supplies
\src{tm215_halts_iff_ZFC_inconsistent} and
\src{tm215_nonhalting_iff_ZFC_consistent}.
Here \src{BProv} means provability from some finite list of
axioms in the given axiom predicate; its letter B does not impose
a fixed proof-length bound. The underlying presentation has the
ZF schemas and the stated choice principle.
The ordinary calibration proof in the repository expressly gives
the implication and converse as reasoning formalizable in ZFC
\cite{kproof}.
Sections~\ref{sec:polynomial}--\ref{sec:complete} give the
mathematical equivalence
\[
 \exists\mathbf x\in\Pos^{15}\ \Pbase(\mathbf x)=0
 \quad\longleftrightarrow\quad
 \Halt(\mathrm{K215},\mathrm{blank}). \tag{77}
\]
Every definition is finite or an elementary recursively defined
integer sequence. The proofs use induction, integer arithmetic,
finite combinatorics, Pell theory, and explicit finite computation
encodings. These constructions and arguments formalize in ZFC.
Combining them with the calibrated K215 theorem proves
Theorem~\ref{thm:main}. This does not infer an internal ZFC theorem
merely from the existence of a Lean theorem: the mathematical
argument supplies the formalizable proof, and the existing
formalization is an independent calibration artifact.

\paragraph{Formal verification status.}
The archive records source identities and the existing endpoint
statements. It does not contain a new Lean/Rocq proof of (77), nor
a kernel-checked encoded ZFC derivation of (1).
The repository's K215 verification dossier also distinguishes
its endpoint theorems from a separately encoded internal ZFC or PA
derivation. We inspected that dossier and the source endpoints;
we did not rerun the entire Lean or Rocq kernel build in this work.
The Python checks certify their stated finite properties, not the
universal mathematical theorem.

\paragraph{Logical consequence.}
No consistency assumption is needed to prove the biconditional.
To conclude that the equation actually has no solution, one must
assume or establish \(\Con(\mathrm{ZFC})\).
Under the usual hypotheses for G\"odel's second incompleteness
theorem, if ZFC is consistent, ZFC cannot prove this absence of
solutions. Mere consistency does not by itself justify the
additional assertion that ZFC cannot prove the opposite
\(\Sigma_1\) statement. A suitable soundness assumption, such as
\(\Sigma_1\)-soundness, suffices for that further conclusion.

\section{Exact degrees and arithmetic counts}\label{sec:degrees}
The pinned source
\pathref{complete84_scaled_strong_output.json} yields the emitted
67-gate circuit by deleting the
ordinary-input cone and its factor, removing the input-dependent
split of \(\gamma\), replacing the marked-word coordinate by the
blank-word \(C=q-F-\alpha\), and retaining the shared packing
expression \(G\). Its complete live ledger is
\(67=39M+28A\). The relation to this parent is source provenance;
soundness was proved above directly for the modified local language,
not by erasing hypotheses from the parent's universal theorem.
Give every unknown degree one and every fixed coefficient degree
zero. Put
\[
 Q=mJ,\quad k_0=\eta+\zeta,\quad C_1=Q-F-\alpha,\quad
 T_0=wC_1-zQ. \tag{78}
\]
Then
\[
 X_{\rm top}=wQ,\quad Y_{\rm top}=sQ^3,\quad
 a_{\rm top}=wsQ^4,\quad c_{\rm top}=k_0sQ^3. \tag{79}
\]
The exact factor degrees in (6) are
\[
 22,\quad18,\quad60,\quad7,\quad2,\quad46. \tag{80}
\]
The main norm has a significant cancellation:
\[
 N_1=2ac(X+\gamma H)+(X+\gamma H)^2-Hc^2. \tag{81}
\]
Its degree-\(18\) leader is \(8\gamma a_{\rm top}^2c_{\rm top}\).
The product's leading homogeneous form is
\[
 -8Q^{91}h\gamma i^4k_0^{13}w^{13}s^{26}T^2f^2T_0. \tag{82}
\]
For example, the degree-\(155\) monomial
\[
 J^{92}h\gamma i^4\eta^{13}w^{13}s^{26}T^2f^2z \tag{83}
\]
has coefficient \(8m^{92}\ne0\). The final subtraction of
\(\Delta\), of degree \(12\), cannot cancel it. The degree is
therefore exactly \(155\), on this fixed coefficient slice.
This calculation does not simplify using identities that hold
only at solutions.
For the reduced auxiliary coefficient,
\[
 \Plow-\Pbase=
 (N_0N_1N_KN_T)N_S(N_S-\Delta)(V^2-y^2). \tag{84}
\]
At any zero of either polynomial with \(\Delta\ne0\),
the integer unit \(N_s\) in (31) must be \(1\), since
\(\Delta\equiv0,3\pmod4\) excludes \(-1\).
Hence \(N_S=\Delta\) and \(S=\Delta(f^2-1)\);
the two auxiliary factors agree. If \(\Delta=0\), both
polynomials vanish identically because \(N_S=0\).
Thus this deformation actually preserves their full integer zero
sets as polynomials in the supplied coordinates, and in particular
their required positive zero sets.
The degree of \(N_A'\) is \(28\), and the new full leader is
\[
 -8Q^{71}h\gamma i^2k_0^9w^{11}s^{20}T^2f^4T_0. \tag{85}
\]
Its \(J^{72}h\gamma i^2\eta^9w^{11}s^{20}T^2f^4z\)
coefficient is \(8m^{72}\ne0\), proving exact degree \(123\).
One extra subtraction is needed, for \(68=39M+29A\).
Deleting the five product-combination gates and final subtraction
from (7) gives the 61-operation system. The analogous system
for (9) costs 62. For a single lower-degree polynomial, sum the
squares of the six residuals in (8). The six residual
subtractions, six squares and five additions add 17 operations.
The maximum residual degrees are \(60\) and \(46\);
their unique highest-degree squares give exact degrees \(120\)
and \(92\). The JSON system packets include both the comparison
interface and these fully charged sum-of-squares finalizers.

\section{An equation over all integers}\label{sec:integers}
The phrase \emph{Diophantine equation} often prescribes variables in
\(\Z\), rather than \(\Pos\). The following explicit conversion
resolves that domain issue without division.
\begin{lemma}[Positive integers with three signed coordinates]
\[
 \Pos=\{1+u^2+v^2+r(r+1):u,v,r\in\Z\}. \tag{86}
\]
\end{lemma}
\begin{proof}
The expression is positive because \(r(r+1)\ge0\).
Conversely let \(n\ge1\). The number \(4(n-1)+1\) is
\(1\pmod4\) and is not of the excluded form
\(4^a(8b+7)\). By Legendre's three-square theorem it is a sum of
three integer squares. Exactly one summand has an odd root.
Write the roots, after reordering, as \(2u,2v,2r+1\)
and divide the equality by four. This proves (86).
The same integer parametrization has a formalized proof in
the Archive of Formal Proofs \cite{threesquares}.
\end{proof}
The coordinates \(\tau,y\) occur only through their squares.
With the other thirteen coordinates positive, neither can be
zero at a zero of \(\Pbase\) or \(\Plow\).
Indeed \(\tau=0\) would make \(N_0<-1\);
after the strong unit is recovered, \(y=0\) would make the
auxiliary unit \(SV^2=1\) with \(S>1\), impossible.
Allowing them to be signed therefore changes no solvability
statement: take absolute values.
For each of the thirteen coordinates in
\[
 \mathcal I=\{J,F,\alpha,z,f,h,i,T,s,w,\eta,\zeta,\gamma\}, \tag{87}
\]
introduce three separate integer unknowns \(u_x,v_x,r_x\), and set
\[
 x=1+u_x^2+v_x^2+r_x(r_x+1). \tag{88}
\]
Keep \(\tau,y\in\Z\). Equations (4)--(7), (25), and (88)
now specify one ordinary polynomial in \(13\cdot3+2=41\)
unrestricted integer variables. Its absence of integer solutions
is provably equivalent in ZFC to \(\Con(\mathrm{ZFC})\).
The packets with suffix \src{over_Z41} spell out every gate.
Each substitution costs seven operations, for total costs
\(67+13\cdot7=158\) and \(68+13\cdot7=159\).
The exact degrees double to \(310\) and \(246\):
the leading forms (82) and (85) contain neither \(\tau\) nor
\(y\), and replacing the other variables by the disjoint
quadratic forms \(u_x^2+v_x^2+r_x^2\) is injective on the
polynomial ring. Set \(v_x=r_x=0\) to see the injective
monomial map \(x\mapsto u_x^2\).

\section{Reproduction, scope, and further work}\label{sec:reproduce}
\subsection{What the archive contains}
The main files are \src{report.tex} and \src{report.pdf}.
The \src{data/} directory contains the literal polynomial and
system circuits, the full transition list, the finite-expression
coefficient recipe, independent cardinality and local-semantics
receipts, and source manifests.
The \src{code/} directory contains standard-library Python
emitters and auditors.
The \src{proofs/} directory supplies additional readable
arithmetic, tableau, and source-calibration notes.
The \src{upstream/} directory preserves a scoped set of
proof dependencies and calibration sources from the pinned
repository revision, with original relative paths.
The applicable upstream licenses are in \src{licenses/}.
From the extracted archive, run:
\begin{lstlisting}
python3 -B code/verify_all.py --output verification_receipt.json
pdflatex -interaction=nonstopmode -halt-on-error report.tex
pdflatex -interaction=nonstopmode -halt-on-error report.tex
\end{lstlisting}
Python needs no third-party packages. The audit uses a temporary
directory and compares rebuilt receipts with the distributed data.
No network connection, historical research script execution,
actual expansion of \(B\), or actual halting witness is needed.
The root \src{README.md} gives the supported options and
the final recorded result.

\subsection{Evidence and its limits}
The finite audits establish:
\begin{itemize}
\item exact agreement of all 430 transition rows in the three
exported K215 representations and the stated hashes;
\item agreement of the local predicate with finite halting
controls, independent padding and helical seam checks, and
rejection of a nonhalting truncation;
\item the exact window cardinality, by closed category counts
and a separate direct predicate enumeration;
\item deterministic reconstruction of all arithmetic circuits,
live ports and gates, operation ledgers, independent exact
evaluations, and degree diagnostics supported by the uniform
leading-form proof;
\item the scalar constant recipe, its no-carry bounds and
modulo-five identity, with independent sparse small-layout
checks of band and anchor geometry.
\end{itemize}
They do not decide whether K215 halts, produce a ZFC refutation,
expand the giant fixed coefficients, search for an actual solution,
or replace the universal mathematical arguments above.
The proof is intended to be reviewable and suitable for subsequent
formalization; the finite checks should not be described as a
kernel certificate of the final consistency equivalence.

\subsection{Useful next improvements}
The most valuable next formal step is to encode the exact
local predicate, native-word soundness, and arithmetic rank
composition in Lean or Rocq and connect them to the existing
K215 endpoint. This would turn the distinction just stated
into a machine-checked bridge.
For mathematical simplification there are several different
objectives. A direct compiler for the overlapping cross relation
could reduce coefficient-description cost substantially by
avoiding the full \(3\)-by-\(3\) window alphabet. It would require
a fresh carry-free and synchronization proof; such a reduction
is not claimed here. Applying an effective nine-variable
reduction to this explicit polynomial is another concrete
project, but its complete specialized formula and coefficient
recipe must actually be emitted before claiming that alternative.
Reducing the remaining positivity witnesses or changing the
Pell scales requires new full-zero-set arguments, since deleting
a coordinate based only on canonical witnesses can create
spurious solutions.
No global simplest Diophantine equation is asserted without
a chosen ordering of variable count, degree, arithmetic cost,
coefficient height, and proof complexity. The archive makes
these tradeoffs measurable for one fully specified consistency
instance.

\appendix
\section{Complete arithmetic circuit}\label{app:gates}
Every row is one binary operation. The two fixed-numeral port
names that need care are \src{Bm1} for \(m\) and \src{MF}
for \(M_F^{\rm src}\). The source register \src{A} denotes
\(\Delta\), not the Pell parameter \(a+2\).
The input aliases are
\src{Jrep}=\(J\), \src{transport_quotient}=\(z\),
\src{auxiliary_quotient}=\(T\), \src{tau_root}=\(\tau\),
and \src{y_aux}=\(y\).
\begingroup\footnotesize
\setlength{\tabcolsep}{3pt}
\begin{longtable}{rll}
\toprule
No. & Output register & Operation \\
\midrule
\endfirsthead
\toprule
No. & Output register & Operation \\
\midrule
\endhead
\bottomrule\endfoot
1 & \src{tau_square} & \src{tau_root * tau_root} \\
2 & \src{repunit} & \src{Bm1 * Jrep} \\
3 & \src{q} & \src{repunit + 1} \\
4 & \src{Lbig} & \src{q * q} \\
5 & \src{n2} & \src{Lbig * q} \\
6 & \src{wn2} & \src{w * q} \\
7 & \src{sn2} & \src{s * n2} \\
8 & \src{UM} & \src{wn2 * sn2} \\
9 & \src{R10b} & \src{eta + zeta} \\
10 & \src{ksn2} & \src{R10b * sn2} \\
11 & \src{first_root_base} & \src{UM * ksn2} \\
12 & \src{first_next} & \src{first_root_base + R10b} \\
13 & \src{first_product} & \src{first_root_base * first_next} \\
14 & \src{norm_first} & \src{tau_square - first_product} \\
15 & \src{R10a} & \src{ksn2 + eta} \\
16 & \src{R12} & \src{UM + sn2} \\
17 & \src{cam2} & \src{R10a * R12} \\
18 & \src{D1} & \src{wn2 + cam2} \\
19 & \src{a4} & \src{4 * R12} \\
20 & \src{a4m5} & \src{a4 + 3} \\
21 & \src{gam} & \src{gamma * a4m5} \\
22 & \src{R14} & \src{D1 + gam} \\
23 & \src{L15} & \src{R14 * R14} \\
24 & \src{a_square} & \src{R12 * R12} \\
25 & \src{A} & \src{a_square + a4m5} \\
26 & \src{c2} & \src{R10a * R10a} \\
27 & \src{Ac2} & \src{A * c2} \\
28 & \src{norm_main} & \src{L15 - Ac2} \\
29 & \src{norm_pair} & \src{norm_first * norm_main} \\
30 & \src{q_minus_F} & \src{q - F} \\
31 & \src{C} & \src{q_minus_F - alpha} \\
32 & \src{aux_y2} & \src{y_aux * y_aux} \\
33 & \src{hpm1} & \src{h * UM} \\
34 & \src{index_difference} & \src{R10b - hpm1} \\
35 & \src{gap_product} & \src{repunit * q_minus_F} \\
36 & \src{gap} & \src{gap_product + alpha} \\
37 & \src{Lm1} & \src{Lbig - 1} \\
38 & \src{rproduct} & \src{gap * Lm1} \\
39 & \src{qMF} & \src{q * MF} \\
40 & \src{mask_factor} & \src{MC + qMF} \\
41 & \src{mask} & \src{mask_factor * Jrep} \\
42 & \src{r_lhs} & \src{rproduct + mask} \\
43 & \src{norm_index} & \src{index_difference - r_lhs} \\
44 & \src{kinner} & \src{Kconstant + w} \\
45 & \src{innerC} & \src{kinner * C} \\
46 & \src{transport_partial} & \src{innerC + q_minus_F} \\
47 & \src{local_rhs} & \src{transport_quotient * repunit} \\
48 & \src{norm_transport} & \src{transport_partial - local_rhs} \\
49 & \src{L16} & \src{f * f} \\
50 & \src{auxiliary_Tf} & \src{auxiliary_quotient * f} \\
51 & \src{auxiliary_Tf_minus_one} & \src{auxiliary_Tf - 1} \\
52 & \src{auxiliary_c_Tf} & \src{R10a * auxiliary_Tf_minus_one} \\
53 & \src{auxiliary_R_f2} & \src{r_lhs * L16} \\
54 & \src{aux_u_rhs} & \src{auxiliary_c_Tf - auxiliary_R_f2} \\
55 & \src{H2} & \src{aux_u_rhs * aux_u_rhs} \\
56 & \src{aux_square_gap} & \src{H2 - aux_y2} \\
57 & \src{aux_coefficient_root} & \src{i * Ac2} \\
58 & \src{R16} & \src{aux_coefficient_root * aux_coefficient_root} \\
59 & \src{scaled_f_square} & \src{A * L16} \\
60 & \src{norm_strong} & \src{scaled_f_square - R16} \\
61 & \src{L17} & \src{R16 * aux_square_gap} \\
62 & \src{norm_aux} & \src{L17 + aux_y2} \\
63 & \src{norm_four} & \src{norm_pair * norm_aux} \\
64 & \src{norm_product} & \src{norm_four * norm_index} \\
65 & \src{all_units} & \src{norm_product * norm_transport} \\
66 & \src{seven_units} & \src{all_units * norm_strong} \\
67 & \src{polynomial} & \src{seven_units - A} \\
\end{longtable}
\endgroup


\section{Independent local-window cardinality}\label{app:counts}
Only the five-position cross enters the local predicate.
The four corners contribute the free factor \(439^4\).
The category table below records the number of valid crosses
per center and the number of center symbols in that category.
The direct checker sums over every left/right pair for a
representative of each category and independently computes its
allowed previous/next combinations. The grouping of moving
head states is valid because the new state of a moving center
head does not occur in that center's own next payload.
\begin{center}\small
\begin{tabular}{lrr}
\toprule
Center category & Multiplicity & Crosses per center \\
\midrule
Vertical boundary & 1 & 84,411,798 \\
Horizontal boundary & 1 & 168,438,154 \\
Headless, symbol 0 & 1 & 184,091,400 \\
Halt, symbol 0 & 1 & 98,097,108 \\
Headless, symbol 1 & 1 & 155,712,504 \\
Halt, symbol 1 & 1 & 98,097,108 \\
Nonhalting, write 0, left & 51 & 99,368,184 \\
Nonhalting, write 0, right & 190 & 99,600,324 \\
Nonhalting, write 1, left & 31 & 71,301,144 \\
Nonhalting, write 1, right & 158 & 71,301,144 \\
Initial L & 1 & 2,118 \\
Initial Q & 1 & 437 \\
Initial R & 1 & 2,110 \\
\midrule
All centers & 439 & Total: 38,256,607,897 \\
\bottomrule
\end{tabular}
\end{center}


\section{Provenance and proof-dependency map}\label{app:provenance}
The source snapshot is
\begin{center}\small
\texttt{cdd582bec13f908df3711fa9e05d1ddac368a5db}.
\end{center}
The immediate arithmetic parent has SHA-256
\begin{center}\small
\texttt{8b2bc5d0b4db90c168107b7ded2cb4ca08df0098ed190851a3db4f1e49a9c0cf}.
\end{center}
The machine hash appears in Section~\ref{sec:machine}.
The machine/source agreement manifest and the complete
archive hash manifest are separate from the polynomial's
mathematical coefficients: hashes identify the supplied bytes,
while (10)--(25) define the mathematical objects.
The imported proof collection is deliberately scoped.
Historical progress headers may point to other research packets;
those headers do not enlarge the dependency claims of this paper.
The detailed machine-readable \src{data/proof_dependencies.json}
records the imported sections and file hashes.
\begin{center}
\begin{tabular}{p{.28\textwidth}p{.64\textwidth}}
\toprule
Obligation & Principal included proofs\\
\midrule
Strong rank and signed target &
Generic relaxed auxiliary rank, signed Pell proof, half-parameter
Pell proof, reversed auxiliary and exact quotient projection\\
Power and binomial extraction &
Base-two Pell proof, half-binomial kernel, asymmetric scale proof\\
Normalized strong converse &
Normalized strong auxiliary witness; explicit divisibility at \(m_1=2cR\)\\
Native masks and alignment &
Fixed native compiler 76; modified half-binomial compiler;
canonical bounded-projection margin\\
Actual-index tuning &
Five-adic Boolean dummy-control theorem\\
Blank fixed-input semantics &
New \src{blank_helical_proof.md} and exact executable predicate\\
ZFC calibration &
K215 ordinary proof, Lean endpoint, Rocq endpoint, verification README\\
\bottomrule
\end{tabular}
\end{center}

\clearpage
\begin{thebibliography}{99}
\raggedright
\bibitem{kendpoint}
V.~Reshetnikov, ProveIt, \emph{K215 Lean endpoint},
\pathref{Computability/BusyBeaver/Lean/BusyBeaver/ZFC/KEndpoint215.lean},
snapshot stated above.
\url{https://github.com/VladimirReshetnikov/ProveIt/blob/cdd582bec13f908df3711fa9e05d1ddac368a5db/Computability/BusyBeaver/Lean/BusyBeaver/ZFC/KEndpoint215.lean}.
\bibitem{kproof}
V.~Reshetnikov, ProveIt, \emph{Candidate K215: ordinary mathematical
calibration proof}, \pathref{Computability/BusyBeaver/Research/ZFC/candidate-k215/proof.md},
and \pathref{k225-formal/verification/k215/README.md}, same snapshot.
\bibitem{jones}
J.~P.~Jones, \emph{Universal Diophantine equation},
Journal of Symbolic Logic \textbf{47} (1982), 549--571.
The original repository's corrected transcription is available at
\pathref{Computability/HilbertTenthProblem/Papers/1982/jones1982_corrected.tex}.
\url{https://doi.org/10.2307/2273588}.
\bibitem{pellrank}
ProveIt, \emph{Relaxed auxiliary Pell rank},
\pathref{Computability/HilbertTenthProblem/Papers/1980/PELL_RELAXED_AUXILIARY_PROOF.md}.
\bibitem{pellsigned}
ProveIt, \emph{Signed Pell proof},
\pathref{Computability/HilbertTenthProblem/Papers/1980/PELL_SIGNED_PROOF.md}.
\bibitem{halfparameter}
ProveIt, \emph{Half-parameter Pell proof},
\pathref{Computability/HilbertTenthProblem/Papers/1980/HALF_PARAMETER_PELL_92_PROOF.md}.
\bibitem{halfbinomial}
ProveIt, \emph{A strong half-binomial Pell kernel in 42 operations},
\pathref{Computability/HilbertTenthProblem/Papers/research-wip/native-stream-queue/pell_kernel_half_binomial42.md}.
\bibitem{asymmetric}
ProveIt, \emph{Asymmetric scale tradeoffs},
\pathref{Computability/HilbertTenthProblem/Papers/research-wip/native-stream-queue/complete75_asymmetric_scale_tradeoffs.md}.
\bibitem{normalized}
ProveIt, \emph{A normalized strong auxiliary witness gives 87 operations},
\pathref{Computability/HilbertTenthProblem/Papers/research-wip/native-stream-queue/complete75_normalized_strong87.md}.
\bibitem{quotient}
ProveIt, \emph{An auxiliary quotient gives a complete universal
polynomial in 85 operations},
\pathref{Computability/HilbertTenthProblem/Papers/research-wip/native-stream-queue/complete85_auxiliary_bezout_projection.md}.
\bibitem{compiler76}
ProveIt, \emph{Fixed native universal compiler},
\pathref{Computability/HilbertTenthProblem/Papers/1980/FIXED_RAW_UNIVERSAL_76_PROOF.md}.
\bibitem{compiler75}
ProveIt, \emph{Modified compiler masks for the half-binomial 75},
\pathref{Computability/HilbertTenthProblem/Papers/research-wip/native-stream-queue/complete75_half_binomial_compiler.md}.
\bibitem{fiveadic}
ProveIt, \emph{Boolean dummy bits can control the packed index
modulo a power of five},
\pathref{Computability/HilbertTenthProblem/Papers/1980/EXPLORATION_FIVE_ADIC_DUMMY_CONTROL.md}.
\bibitem{threesquares}
A.~Danilkin and L.~Chevalier, \emph{Three Squares Theorem},
Archive of Formal Proofs, 3 May 2023.
\url{https://isa-afp.org/entries/Three_Squares.html}.
\end{thebibliography}
\end{document}
